\section*{第 \ref{ch:b2:climate-variability} 章——气候变率}

\begin{solution}{exo:b2:climate-variability:1}
(a) $29.27 - 26.71 = \qty{2.56}{\celsius}$（2026 年）；$28.79 - 26.73 = \qty{2.06}{\celsius}$（1997 年）。
(b) 2026 年 9 月既是较暖的月份，高出 \qty{0.48}{\celsius}，也是距平较大的月份，
高出 \qty{0.50}{\celsius}。(c) 每一年都与以其为中心的 30 年基准期相比较，基准期每
五年推移一次；1980 年代和 1990 年代的 9 月\omterm{def:b2:climate-variability:modes}{气候态}，
与 2010 年代和 2020 年代的略有不同。
\end{solution}

\begin{solution}{exo:b2:climate-variability:2}
(a) JJA：$(0.79 + 0.95 + 1.28)/3 = 1.007$；JAS：$(0.95 + 1.28 + 1.55)/3 = 1.260$；ASO：$(1.28 +
1.55 + 1.68)/3 = 1.503$；SON：$(1.55 + 1.68 + 2.01)/3 = 1.747$；OND：$(1.68 + 2.01 + 2.02)/3 =
1.903$，单位均为 \unit{\celsius}。(b) 四舍五入至一位小数：1.0、1.3、1.5、1.7、1.9，与表中
数值一致。(c) ASO 2023 时为 1.5。
\end{solution}

\begin{solution}{exo:b2:climate-variability:3}
(a) 该指数为塔希提减达尔文的气压距平：塔希提低于平均，
达尔文高于平均，指数为负。(b) \omterm{def:b2:climate-variability:enso}{厄尔尼诺}。(c) 信风减弱，甚至
西部近地面出现西风，印度尼西亚上空降水比往常偏少，因为
深对流已向东移动。
\end{solution}

\begin{solution}{exo:b2:climate-variability:4}
(a) $804.1/880.6 = \qty{91.3}{\percent}$，为平均值的比例，处于 90 到 \qty{110}{\percent} 之间，
属正常季节，但位于下段。(b) 季节总量为 $556 + 659 + 427 + 252 =
\qty{1894}{mm}$；7 月占其中 $659/1894 = \qty{34.8}{\percent}$，超过三分之一。
\end{solution}

\begin{solution}{exo:b2:climate-variability:5}
(a) 对 \qty{50}{m}，$\tau = 1025 \times 3992 \times 50/16.1 = \qty{1.271e7}{s} = \qty{147.1}{\day}$；
对 \qty{100}{m} 时为其两倍，即 \qty{294.2}{\day}。(b) $\mathrm{e}^{-30.44/147.1} =
0.813$，$\mathrm{e}^{-30.44/294.2} = 0.902$。(c) 阻尼时间与 $h$ 成正比：
混合层越浅，相同通量下储存的热量越少，距平衰减越快，因此
海洋的“记忆”及其所带来的可预报性随之下降。
\end{solution}

\begin{solution}{exo:b2:climate-variability:6}
(a) 三个月为 91.3 天：$1.2 \times \mathrm{e}^{-91.3/147.1} = \qty{0.65}{K}$；六个月
为 182.6 天：$1.2 \times \mathrm{e}^{-182.6/147.1} = \qty{0.35}{K}$。(b) 距平衰减
较慢，因此它作为下月距平的估计优于零。但随机强迫每天都会
产生新的距平，这些距平在预报时未知；随着初始距平衰减，
它们所占比重越来越大，预报逐渐趋向\omterm{def:b2:climate-variability:modes}{气候态}，在那里
预报已无帮助。
\end{solution}

\begin{solution}{exo:b2:climate-variability:7}
(a) $160 \times 86\,400 = \qty{1.38e7}{J}$，即每平方米每天 \qty{13.8}{MJ}。(b) 该
层每平方米含空气 $(1000 - 850) \times 100/9.81 = \qty{1529}{kg}$，热容为
$1529 \times 1004.5 = \qty{1.536e6}{J.m^{-2}.K^{-1}}$；若全部吸收，则每天升温 $1.38 \times
10^7/1.536 \times 10^6 = \qty{9.0}{K}$。(c) 该层以红外线形式向
太空和上方空气辐射热量；暖空气上升，与上方气层混合，或
被风带走，其中包括热低压自身吸入的较冷空气。
\end{solution}

\begin{solution}{exo:b2:climate-variability:8}
(a) $(-1.93 - 1.11 - 1.98)/3 = -1.673$，约为 $-1.67$。(b) 这是有记录以来最低的冬季。
(c) \omterm{def:b2:climate-variability:nao}{北大西洋涛动}处于强负位相：冰岛低压与亚速尔高压偏弱，急流减弱，风暴
路径比往常偏南；北欧为寒冷干燥的冬季，南欧为湿润的冬季，
美国东部冷空气爆发增多。
\end{solution}

\begin{solution}{exo:b2:climate-variability:9}
(a) $9.3 - 0.4 = \qty{8.9}{mm}$ 每天，超过 2.5，在该区域内。(b) 在南半球
夏季为 12 月–次年 2 月：$6.1 - 4.2 = \qty{1.9}{mm}$ 每天，小于 2.5，不在该区域内，
虽然该处降水多于区域内大多数地方。(c) $1894/122 = \qty{15.5}{mm}$
每天；喀拉拉的冬季降水速率只有超过 $15.5 - 2.5 = \qty{13.0}{mm}$ 每天，才会位于区域之外。
(d) \omterm{def:b2:climate-variability:monsoon}{季风}是风向的季节性反转：像第二个
站点那样全年湿润的地方没有\omterm{def:b2:climate-variability:monsoon}{季风}，而雨季短促而强烈的干燥地区则有\omterm{def:b2:climate-variability:monsoon}{季风}。
\end{solution}

\begin{solution}{exo:b2:climate-variability:10}
(a) 求导得 $\ddot T = a\dot H = -abT$：此为角
频率 $\omega = \sqrt{ab}$ 的谐振子方程，周期为 $2\pi/\sqrt{ab}$。(b) 由 $2\pi/\sqrt{ab} = 4$ 得 $ab =
(2\pi/4)^2 = \qty{2.47}{\text{年}^{-2}}$。(c) 令 $T = u\,\mathrm{e}^{-rt}$，$H = v\,\mathrm{e}^{-rt}$，
$\dot T = (\dot u - ru)\,\mathrm{e}^{-rt} = (av - ru)\,\mathrm{e}^{-rt}$，于是 $\dot u = av$，且
同理 $\dot v = -bu$：$(u, v)$ 做无阻尼振荡，而 $(T, H)$ 即为该振荡
乘以 $\mathrm{e}^{-rt}$。一个周期内振幅降为原来的 $\mathrm{e}^{-0.25 \times
4} = \mathrm{e}^{-1} = 0.37$。(d) 强迫之一是信风的随机涨落，它
反复扰动系统；另一是\omterm{def:b2:climate-variability:bjerknes}{皮耶克尼斯反馈}，它使阻尼减弱或
甚至变为负值。
\end{solution}

\begin{solution}{exo:b2:climate-variability:11}
(a) $A = R^2 \,\Delta\lambda\,(\sin\ang{5} - \sin(\ang{-5})) = (6.371 \times 10^6)^2 \times 0.8727
\times 0.1743 = \qty{6.17e12}{m^2}$. (b) $Q = 1025 \times 3992 \times 100 \times 2 \times 6.17
\times 10^{12} = \qty{5.05e21}{J}$。(c) 六个月为 $182.6 \times 86\,400 = \qty{1.578e7}{s}$；
每平方米，$8.18 \times 10^8/1.578 \times 10^7 = \qty{51.9}{W/m^2}$，持续半
年。(d) 来自海洋自身：信风减弱，东部上翻的冷水减少，而
西部储存在较深温跃层之上的暖水向东扩展；变暖的大部分
可能来自海洋中原有的热量，是被搬运而非新增的。
\end{solution}

\begin{solution}{exo:b2:climate-variability:12}
(a) 回路返回 $0.3g$，继而为 $0.3g^2$，依此类推：总和为 $0.3(1 + g + g^2 + \dots) =
0.3/(1 - g)$，这是当 $g < 1$ 时收敛的等比级数：\qty{0.6}{K}、\qty{1.5}{K} 及
\qty{3.0}{K} 分别对应 $g = 0.5$、0.8 与 0.9。(b) 该级数发散：距平将无限增长，
无界；增长须由回路未包括的过程制止，阻尼则来自地表
热通量，以及充放振子所描述的暖水储库的泄放。(c)
海表越暖，微小的增温就越容易引发深对流和降水，
其加热驱动风场；这正是 ENSO 影响最强的季节。
\end{solution}

\begin{solution}{pb:b2:climate-variability:1}
\textbf{1.} $(0.26 + 0.75 + 1.13)/3 = \qty{0.713}{\celsius}$.
\textbf{2.} MJJ $(0.75 + 1.13 + 1.52)/3 = 1.133$；JJA $(1.13 + 1.52 + 1.80)/3 = 1.483$；JAS $(1.52 +
1.80 + 2.06)/3 = 1.793$.
\textbf{3.} ASO $(1.80 + 2.06 + 2.26)/3 = 2.040$；SON $(2.06 + 2.26 + 2.42)/3 = 2.247$；OND $(2.26
+ 2.42 + 2.36)/3 = 2.347$；NDJ $(2.42 + 2.36 + 2.33)/3 = 2.370$。
\textbf{4.} 四舍五入到一位小数：0.7、1.1、1.5、1.8、2.0、2.2、2.3 和 2.4，即表中数值。
\textbf{5.} 单月含有天气噪声；ENSO 持续许多个月，三个月的平均
可去除大部分噪声，同时保留缓慢的信号。
\textbf{6.} 1997 年 JJA 为 1.483，四舍五入为 1.5；其在 NDJ 达到峰值 2.4。
\textbf{7.} 自 2018 年 ASO 至 2019 年 MJJ，以及自 2019 年 SON 至 2020 年 FMA。
\textbf{8.} 十个季节，以及六个季节。
\textbf{9.} 2020 年 ASO；2020 年内仅有四个季节，即 ASO 至 NDJ，不高于 $-0.5$，
差一个才满足规则，而 2021 年 DJF 补足了它。
\textbf{10.} 一次\omterm{def:b2:climate-variability:enso}{厄尔尼诺}，自 2018 年 SON 至 2019 年 MAM，共七个季节（AMJ 2019 降至 0.37）；
\omterm{def:b2:climate-variability:enso}{拉尼娜}始于 2020 年 AMJ（$-0.55$），比 ONI 所定的 ASO 早四个季节。
\textbf{11.} 否：2019 年末相对指数从未超过 0.29。当时整个热带海洋
比其基准期偏暖约半摄氏度；Niño-3.4 区域按绝对值计偏暖，
但几乎不比周围热带地区更暖。
\textbf{12.} 十一。
\textbf{13.} 8、7、10、4、5、6、12、5、9 和 3 年。
\textbf{14.} 最短为 3 年（2023 至 2026 年），最长为 12 年（1997 至 2009 年）。
\textbf{15.} 两种算法结果均为 $69/10 = 6.9$ 年：相继间隔之和即从
第一次开始到最后一次开始的时间，$2026 - 1957 = 69$ 年。
\textbf{16.} 该事件已持续五个季节，即自 MAM 至 JAS，均不低于 \qty{0.5}{\celsius}，并且
峰值为 2.2，高于 1.5；其起始季节 MAM 2026 是该段的第一个季节，之后的季节
只能延长该段，而不能改变其起点。
\textbf{17.} 否，而且也不必如此：全部强度的 25 次事件大约每三年出现一次，正如 FAQ 所说；
强事件只是其中一部分，因此更为罕见，其间隔在 3 至 12 年之间。
\textbf{18.} $\tau = 1025 \times 3992 \times 50/16.1 = \qty{1.27e7}{s} = \qty{147}{\day}$.
\textbf{19.} 混合层：$1025 \times 3992 \times 50 = \qty{2.05e8}{J.m^{-2}.K^{-1}}$；空气柱：
$(10^5/9.81) \times 1004.5 = 10\,194 \times 1004.5 = \qty{1.02e7}{J.m^{-2}.K^{-1}}$，
仅为混合层的 1/20。
\textbf{20.} 海水厚度为 $1.024 \times 10^7/(1025 \times 3992) = \qty{2.5}{m}$。
\textbf{21.} $1.024 \times 10^7/16.1 = \qty{6.36e5}{s} = 7.4$~天。
\textbf{22.} 仅空气的距平约在一周内衰减，混合层的
距平约需五个月，而温跃层之上的暖层更深：只有海洋能
把热量保存足够长的时间，将距平带到下一个季节，而大气提供
使二者耦合的风。
\textbf{23.} 自 1950 年以来，强\omterm{def:b2:climate-variability:enso}{厄尔尼诺}事件平均每隔 \textbf{6.9 年}开始一次，
约七年。
\end{solution}
